\subsection{Limitations}
\label{subsec:limitations}

Several limitations qualify this survey. First, the benchmark corpus
analyzed in Section~\ref{sec:evaluation_benchmarking} is a 2024--2026 snapshot;
fifteen of the eighteen benchmarks were arXiv preprints at the time of
writing and may have been peer-reviewed or superseded since. Claims
that depend on this corpus should therefore be read as time-bounded.
Second, the capability coding behind Figure~\ref{fig:benchmark_heatmap} was
performed by a single reviewer; we have released the per-benchmark
notes as supplementary material to enable community validation, but we
do not report formal inter-rater agreement. Third, the coverage
analysis is biased toward English-language and Anglophone-dominated
benchmarks, which under-represents code-switching, dialect, and
non-Latin script regimes. Fourth, the deployment and safety syntheses
draw heavily on system descriptions from authors of those systems,
which limits adversarial scrutiny of reported latency and security
measurements. Fifth, the proposed Evaluation Card and Deployment Card
templates are not yet empirically validated against community use;
they are framing tools for reviewer consensus, not measurement
instruments. We acknowledge these limitations to preempt
over-generalization of the lifecycle audit and to invite community
extension of the supplementary material.

\section{Conclusion}
\label{sec:conclusion}

This survey reviewed Speech Large Language Models by connecting training,
evaluation, deployment, and speech-specific risk analysis. Training
advances such as Whisper-style weakly supervised pretraining, alignment
modules including AlignFormer and Q-Former,
parameter-efficient fine-tuning, and preference alignment via DPO and
RLAIF have made Speech LLMs increasingly capable. Evaluation, however,
remains a bottleneck: current benchmarks over-measure inherited ASR and
emotion tasks while under-measuring latency, safety, code-switching, and
privacy. Deployment research has begun to address streaming, edge, and
full-duplex constraints, yet
integrated systems combining serving, architecture, and algorithmic
optimization remain rare. Safety, privacy, and governance work has shown
that text-only safeguards do not transfer to speech: audio jailbreaks,
voice cloning, deepfake misuse, and speaker re-identification require
modality-specific responses. By treating Speech LLMs as systems rather
than only models, the lifecycle perspective surfaces both the structural
gaps and the cross-stage couplings that will shape the next generation
of robust, efficient, and responsible spoken-language systems.
